% ==========================================================
% 多个禁位条件（插空法与容斥原理）
% 难度：★★ 中档
% ==========================================================

\newcommand{\qProbPermMultiForbidStem}{%
七名队员站成一横排，要求：甲不站最左端；乙不站最右端；丙不站正中间。求理论上的排法数。
}

\newcommand{\qProbPermMultiForbidAnswer}{%
$3216$
}

\newcommand{\qProbPermMultiForbidSolution}{%
总排列数为 $7!$。

设事件 $E_1$：甲在最左端；$E_2$：乙在最右端；$E_3$：丙在正中间。\\
每一个禁位事件发生时，\\
固定一人，其余六人排列，方法数为 $6!$。\\
任意两个禁位事件同时发生时，\\
固定两人，其余五人排列，方法数为 $5!$。\\
三个禁位事件同时发生时，\\
三人位置固定，其余四人排列，方法数为 $4!$。

由容斥原理，合法方法数 $N$ 等于：\\
总排列数 - 至少违反一个条件的情况数。
\[
N = 7! - (3\times 6!) + (3\times 5!) - 4!
\]
\[
N = 5040 - 2160 + 360 - 24 = 3216
\]
故理论上的排法数为 $3216$ 种。
}

\newcommand{\qProbPermMultiForbidTeachingNotes}{%
\textbf{【避坑与边界说明】}
\begin{itemize}
    \item 七人排列的正中间是第 $4$ 个位置。
    \item 三个禁位彼此不独立，不能直接乘概率。
    \item 也可尝试先排其余四人再插空甲乙丙。\\
    但追踪端点和中心变化太繁琐，容斥更稳。
    \item 容斥公式：两两交集减回(加上)，\\
    三者交集加回(减去)。
    \[ |A\cup B\cup C| = \sum|A| - \sum|A\cap B| + |A\cap B\cap C| \]
\end{itemize}
}
